% !TEX root = ../main.tex
\section*{Chapter 2: Background}
\setcounter{section}{2} \setcounter{subsection}{0}
\phantomsection \label{chap:background} \addcontentsline{toc}{section}{\textit{Chapter 2: Background}}



\subsection{Agent Systems: Prior Work}
\label{sec:agent-systems}

Modern operating systems are built around the process as the primary unit of protection: address-space isolation, controlled file access, and well-defined IPC \cite{tanenbaum2022}.
These abstractions remain fundamental.
Agent-based workloads introduce a layer of intent between the user and execution that they do not directly address: an agent plans over semantically described tools, constructs structured requests, and acts on behalf of a user.
Correctness therefore depends on process isolation and on whether the requested semantics match the operation that actually runs.

Several lines of work have responded to this gap.
Systems such as generative agents demonstrate long-lived entities capable of maintaining memory and coordinating workflows~\cite{park2023generative}.
Frameworks including LangChain~\cite{langchain2022}, LangGraph~\cite{langgraph2024}, AutoGPT~\cite{autogpt2023}, CrewAI~\cite{crewai2024}, and Semantic Kernel~\cite{semantickernel2023} structure tool composition in multi-step pipelines.
More ambitious proposals embed LLM-centric services directly into a system-level kernel managing scheduling, memory, and tool access~\cite{packer2023memgpt, wu2024oscopilot, mei2024aios, liu2026agentos}.
At the application layer, OpenClaw~\cite{openclaw2026} acts as a persistent user-space control layer for tool execution, while OpenHarness~\cite{openharness2026} and Hermes Agent \cite{nousresearch2026hermesagent} integrate control logic directly into the agent process itself; Harness~\cite{harness2025} embeds governed AI agents in CI/CD pipelines; Apple Intelligence~\cite{apple2024pcc} integrates agents with on-device OS services and a private-cloud compute layer.

Section~\ref{sec:limits-user-space} details why this enforcement placement matters.



\subsection{What MCP Provides}

The Model Context Protocol (MCP), released by Anthropic in November 2024, provides a useful abstraction for structuring tool-mediated agent interaction~\cite{mcp2025}.
It defines a client-host-server architecture in which tools, resources, and prompts are exposed through a standardised JSON-RPC interface.
The host coordinates sessions and mediates access, while servers implement capabilities; tools are the primary mechanism by which agents take executable actions on external systems.
MCP rapidly achieved broad ecosystem adoption, with implementations in Python, TypeScript, C\#, and Java, and integration across major LLM providers and agent frameworks~\cite{mcp_anniversary_2025}.

This separation simplifies interoperability and modularity.
However, it also concentrates responsibility in the host: identity management, approval decisions, and consistency between tool descriptions and execution all depend on correct host behavior.
MCP defines how components communicate, but does not specify how a local system should enforce these properties.
linux-mcp implements the enforcement layer that MCP leaves unspecified: it provides the local mechanism that turns the protocol's abstract trust model into a set of kernel-verified invariants.



\subsection{Limits of User-space Mediation}
\label{sec:limits-user-space}

\textbf{Semantic drift.}
Tool invocation depends on semantically rich manifests combining natural-language descriptions, schemas, and examples \cite{schick2023toolformer, patil2023gorilla, mcp2025}.
Because both interpretation and enforcement reside in a mutable user-space process, a change in implementation or a weakened control path may result in executing a capability that no longer matches the approved semantics \cite{liu2023prompt}.

\textbf{Identity anchoring.}
Agent systems bind requests to sessions or agent identifiers, not PIDs \cite{yao2023react, autogpt2023, park2023generative}, and MCP provides no mechanism to anchor these identities in OS-visible credentials.
Prompt injection and session confusion attacks can therefore cause agents to exceed their intended authority even when the gateway appears to enforce correct behaviour \cite{greshake2023prompt, openclaw_security1, openclaw_security2}.

\textbf{Durability.}
Approval decisions and pending requests may be lost on crash or restart when control-plane state is held only in memory \cite{packer2023memgpt}.
The fragility is particularly acute in long-running agent systems where execution spans multiple steps and decisions \cite{wang2023voyager, yang2024swe}.

These limitations are not unique to any single framework; they reflect a broader pattern in which control-plane invariants \cite{lim2024compartmentalization} are difficult to enforce when confined to user space \cite{schneider2022sok, lefeuvre2024sok, ji2026seagent, shi2025progent}.



\subsection{Relation to Native Linux Security Mechanisms}

Any such design must coexist with existing Linux security mechanisms.
Facilities such as seccomp, AppArmor, SELinux, and Landlock provide strong guarantees over system calls, resources, and process privileges \cite{seccomp2026, lsm2026, apparmor2026, selinux2026, landlock2026}.
These mechanisms operate at the level of kernel objects and access rights.

However, none of them capture higher-level concepts such as tool identity, semantic equivalence between a manifest and an execution endpoint, or the scope of an approval decision.
The root cause is a layer mismatch: existing mechanisms enforce policy at the level of kernel objects (files, processes, syscalls), whereas tool identity is an application-layer concept defined by a semantic hash of a JSON manifest.
Extending these mechanisms to enforce tool-identity would require either encoding tool hashes into LSM policy or hooking into the application-protocol layer, which breaks the principle that LSM hooks attach to kernel operations, not to application messages.

This gap motivates introducing a dedicated control layer that validates semantic and identity invariants at the point of tool execution, outside the mutable execution path of the user-space gateway.


\begin{table}[t]
\centering
\footnotesize
\setlength{\tabcolsep}{5pt}
\renewcommand{\arraystretch}{1.25}
\begin{tabularx}{\textwidth}{@{}l >{\columncolor{gray!8}}X X X@{}}
\toprule
\textbf{Property}
  & \textbf{linux-mcp (this work)}
  & \textbf{OpenClaw~\cite{openclaw2026}}
  & \textbf{Hermes Agent~\cite{nousresearch2026hermesagent}} \\
\midrule
Control-plane placement
  & Split: manifest logic in user space; admission and tickets in kernel module
  & Standalone user-space gateway
  & In-process: same PID as the agent loop \\
Trust boundary of admission
  & Kernel-enforced, gated by \texttt{CAP\_NET\_ADMIN}
  & Gateway process only
  & Python runtime; optional Docker sandbox \\
Approval state
  & Kernel-resident tickets with explicit state machine
  & In-memory gateway state
  & Per-session module state in the agent process \\
Session $\leftrightarrow$ identity binding
  & Hash over \texttt{SO\_PEERCRED} re-checked every request
  & Gateway-managed session table
  & Session-id token; no OS-level binding \\
State on gateway crash
  & Registry, tickets, audit log persist in kernel
  & Lost with the gateway process
  & Lost with the agent process \\
Audit availability post-crash
  & Readable through sysfs without a running daemon
  & Application log only
  & Application log only \\
\bottomrule
\end{tabularx}
\caption{Control-plane properties of linux-mcp compared with two representative user-space tool-invocation frameworks, as configured in their default upstream documentation without external sandboxing layers. Production deployments commonly stack additional isolation on top: OpenClaw is often paired with NVIDIA NemoClaw and OpenShell sandboxing, and Hermes Agent's gateway architecture also offers Docker, SSH, Daytona, Singularity, and Modal terminal backends with container hardening. Such deployments would shift the right two columns toward stronger admission and execution boundaries; the comparison here is against the bare framework. The grey column marks the system proposed in this work.}
\label{tab:comparison}
\end{table}


\subsection{Related Work}
\label{sec:related-work}

Suppose an attacker replaces a registered tool's JSON-RPC payload at runtime, so that the \texttt{operation} field now targets a different backend handler.
The socket address and the enclosing syscalls stay identical.
An LSM hook at \texttt{security\_\allowbreak socket\_\allowbreak recvmsg} fires on the same kernel operation as before; a seccomp filter sees the same syscall numbers; a Landlock policy sees the same file descriptors.
None of them has a place to hang a check that says ``this tool's manifest hash no longer matches what was registered''.
That is the gap this section is about.

Existing kernel-native mechanisms such as seccomp-BPF, eBPF/LSM hooks, Landlock, and capability-based designs provide strong control over syscalls, kernel objects, and process-scoped authority~\cite{seccomp2026, wright2002lsm, landlock2026, salaun2017landlock, watson2010capsicum, brimhall2023mac, lefeuvre2024sok}.
They operate below the application-protocol layer, so they do not natively express MCP-specific concepts such as manifest-derived tool identity or semantic equivalence between a declared tool and an executed endpoint.
A tool's identity in this system is a SHA-256 digest of its JSON manifest, a concept that does not exist at the syscall or VFS layer; no hook in seccomp or LSM fires at the point where a JSON-RPC payload is interpreted.
Modern eBPF can perform bounded map lookups and keep per-session state in pinned maps, so neither hash checks nor ticket storage is the hard part. 
The harder part lies at the attachment point. 
eBPF LSM hooks fire on kernel operations such as socket receive and file open, not on parsed JSON-RPC messages, so any in-kernel enforcer would have to re-parse MCP framing and reconstruct tool identity from raw socket bytes on every hook invocation, all under the eBPF verifier's bounded-loop and instruction-count constraints.
For this reason, these mechanisms can harden the surrounding execution environment, but they do not directly enforce the control-plane invariants targeted in this work.

The agent frameworks surveyed in Section~\ref{sec:agent-systems} locate the full control plane in user space; correctness depends entirely on the in-process runtime remaining uncompromised.
OpenClaw~\cite{openclaw2026} operationalises this pattern at scale, and independent security analyses confirm the resulting structural vulnerabilities~\cite{openclaw_security1, openclaw_security2}.
Industrial deployments such as Harness~\cite{harness2025} add platform-level Role-Based Access Control (RBAC) on top, but the trust anchor simply moves to a cloud control plane; it is not eliminated.
Recent systems on application-layer mandatory access control including privilege-escalation defences for LLM agents~\cite{ji2026seagent} and programmable privilege control for tool-calling agents~\cite{shi2025progent} have enforced policy inside a modified agent runtime, so invariant correctness depends on the runtime itself remaining intact.
Most of these approaches keep the trust anchor in user-space or in a cloud-hosted control plane; their correctness degrades when the runtime or gateway is compromised.
Apple Intelligence is a partial exception: its Private Cloud Compute design relies on Apple-controlled hardware, a hardware root of trust, and attestation mechanisms, rather than a general-purpose Linux host control plane.~\cite{apple2024pcc}.
Table~\ref{tab:comparison} contextualises linux-mcp against two representative user-space frameworks, the monolithic-gateway OpenClaw~\cite{openclaw2026} and the in-process Hermes Agent~\cite{nousresearch2026hermesagent}, across the control-plane properties most relevant to this work.


AIOS-style proposals such as AIOS, OS-Copilot, MemGPT, and AgentOS~\cite{mei2024aios, liu2026agentos, packer2023memgpt, wu2024oscopilot} pursue a broader scope: they embed LLM-specific abstractions (scheduling, memory management, tool access) into a new OS substrate.
These proposals broaden scope by redesigning OS infrastructure; linux-mcp instead adds a narrow kernel role to an otherwise unchanged Linux system, targeting a bounded set of MCP-style control-plane invariants.


\begin{table}[t]
\centering
\footnotesize
\setlength{\tabcolsep}{5pt}\renewcommand{\arraystretch}{1.15}
\begin{tabularx}{\linewidth}{@{}l X X@{}}
\toprule
\textbf{Attack class} & \textbf{Representative attack} & \textbf{Defence in linux-mcp} \\
\midrule

Spoofing
& Forged \texttt{session\_id} (including same-UID hijack) or unregistered \texttt{tool\_hash}.
& \texttt{SO\_PEERCRED} binding hash over $(\text{uid},\text{gid},\text{pid},\text{session\_id})$ checked on every request; \texttt{tool\_hash} must match a kernel-resident manifest entry. \\

Replay
& Ticket reused past its TTL, or after the referenced tool was re-registered.
& Tickets carry \texttt{(binding\_hash, binding\_epoch)} + TTL; re-registering a tool bumps its catalog epoch, invalidating dependent tickets and forcing session rebind. \\

Substitution
& Manifest edited post-registration, or backend binary swapped behind a live socket.
& Kernel-verified manifest SHA-256; first successful probe pins a TOFU \texttt{binary\_hash} over \texttt{/proc/<pid>/exe} (composite for interpreters), re-checked on every invocation. \\

Escalation
& Tool invoked without a valid approval via an \texttt{mcpd} logic bug (missed state transition, mis-routed approval, ticket-state TOCTOU).
& Kernel is the authoritative store for ticket state, so forwarding errors in \texttt{mcpd} surface as kernel-side rejections instead of silent admissions. \\

Audit loss
& Kill \texttt{mcpd} to erase recent call history.
& Per-agent \texttt{call\_log} is a fixed-size kernel ring buffer, exposed read-only via sysfs; survives daemon termination. \\

\bottomrule
\end{tabularx}
\caption{Control-plane threats and defences (in scope). The last row is an observability property, not a classical attack class.}
\label{tab:threat_in_scope}
\end{table}


\subsection{Threat Model}
\label{sec:threat-model}

The above limitations suggest that the correctness of MCP-style systems depends on a small set of control-plane invariants, including identity binding, approval validity, and consistency between tool semantics and execution.
We therefore ground our design in a threat model that focuses on attacks against these invariants in a local-host setting.

The local-adversary model is appropriate for MCP-style deployments for three reasons.
First, typical MCP deployments co-locate tool services with the gateway on a single host, so the control plane itself is a local attack surface.
Second, \texttt{mcpd} exposes no network interface but only a Unix domain socket, so a remote adversary without on-host code execution cannot reach the control plane at all, while one with code execution is indistinguishable from a local unprivileged process; the local model therefore captures both cases without a separate remote treatment.
Third, two threats frequently raised in the AI agent setting fall outside this model because they live at other layers, namely prompt injection at the model's reasoning layer and supply-chain compromise of manifests at the filesystem integrity layer; neither is an appropriate target for kernel-level admission control, and both are addressed by complementary defences at those layers, as discussed later in this section.
The threat model therefore focuses on what kernel enforcement can actually address: a co-located unprivileged process that can observe, forge, or replay control-plane messages exchanged with the gateway.

Concretely, the adversary model grants the attacker four capabilities: (1) arbitrary code execution under an unprivileged uid distinct from the one under which \texttt{mcpd} runs; (2) access to any Unix domain socket the gateway exposes to unprivileged processes; (3) the ability to craft and replay arbitrary JSON-RPC messages on that socket, including messages observed in prior interactions; and (4) full knowledge of the MCP protocol, tool manifests, and linux-mcp implementation.
Correspondingly, the attacker interacts with linux-mcp exclusively through the \texttt{mcpd} JSON-RPC surface and holds none of the three elevated capabilities that would bypass this interface: no \texttt{CAP\_\allowbreak NET\_\allowbreak ADMIN}, no \texttt{ptrace} access to \texttt{mcpd}, and no write access to the manifest directory.
Each boundary closes a distinct bypass path: \texttt{CAP\_\allowbreak NET\_\allowbreak ADMIN} would let the attacker talk to the kernel module directly via Netlink, skipping \texttt{mcpd} entirely; \texttt{ptrace} would let the attacker read or modify \texttt{mcpd}'s address space and forge control-plane state from within the gateway; and manifest write would poison the kernel registry at the next registration, since manifests are the trusted input from which tool identity is derived.

Two dominant attack categories in AI agent systems fall outside this scope.
\emph{Prompt injection} operates at the model's reasoning layer, before any \texttt{tool:exec} request is submitted; kernel admission control has no visibility into the planning process \cite{greshake2023prompt, liu2023prompt}.
\emph{Supply-chain compromise of manifests} requires file-system-level defences such as Integrity Measurement Architecture or device-mapper verity; linux-mcp treats on-disk manifests as trusted input and assumes the adversary lacks write access to \texttt{tool-app/\allowbreak manifests/}.

Trust in \texttt{mcpd} is structural, not continuous.
The daemon is trusted to perform two translations correctly: at startup, converting manifest contents into SHA-256 hashes and registering them with the kernel; and at each session open, converting the client's \texttt{SO\_\allowbreak PEERCRED} into a binding hash anchored in the kernel.
Once these registrations and bindings are in place, the kernel holds canonical state and re-checks it on every \texttt{TOOL\_\allowbreak REQUEST} and \texttt{APPROVAL\_\allowbreak DECIDE}, so the daemon's runtime correctness is no longer load-bearing for individual request validation.
This placement defends against (a) daemon logic bugs in request forwarding such as session-id mismatch, missed ticket transitions, or concurrent TOCTOU on ticket state, (b) daemon crashes, since tickets and audit state survive in the kernel, and (c) unprivileged co-located adversaries forging UDS traffic, since session binding is anchored to peer credentials they cannot spoof.
It does \emph{not} defend against a \texttt{mcpd} whose process memory is directly compromised at any point in its lifetime: a compromised daemon still holds a valid \texttt{CAP\_\allowbreak NET\_\allowbreak ADMIN} handle and can drive the Netlink interface with apparently legitimate calls, so runtime compromise of the gateway is explicitly out of scope and is not claimed to be contained by the kernel path.
Narrowing this boundary further would require splitting the \texttt{CAP\_\allowbreak NET\_\allowbreak ADMIN} handle into a minimal separate service that only forwards pre-validated Netlink frames on behalf of \texttt{mcpd}, so that a memory-safety compromise of the main daemon no longer immediately grants direct kernel reach; we return to this in Section~\ref{sec:further-work}.

Within this scope, we focus on four classes of control-plane failures:
(1) \emph{spoofing} of identity or session context,
(2) \emph{replay} of previously valid requests or approval decisions,
(3) \emph{substitution} of tool identity or semantics, and
(4) \emph{escalation} through bypassing admission checks in the gateway.
A fifth property, \emph{post-compromise audit continuity}, sits alongside these four: a gateway crash should not erase the record of recent activity.
Table~\ref{tab:threat_in_scope} summarises the five and the corresponding kernel-anchored defences.
Three of those defences go beyond manifest-hash admission alone: trust-on-first-use pinning of the backend binary identity, per-tool catalog-epoch rebind on manifest change, and a fixed-size kernel call log that survives daemon failure.

The trusted computing base is the kernel module, the operator's TOML configuration, and the manifest directory.
\texttt{mcpd} is trusted to populate the kernel's view of the catalog and to bind sessions at connection time; once populated, the integrity of tickets, tool identity, and audit state does not depend on the daemon staying alive.
The supplied systemd unit (Section~\ref{sec:daemon-implementation}) runs \texttt{mcpd} as an unprivileged user with only \texttt{CAP\_\allowbreak NET\_\allowbreak ADMIN} and \texttt{CAP\_\allowbreak SYS\_\allowbreak PTRACE}; this is defence-in-depth, not part of the model itself.
